\documentclass[12pt,a4paper]{article}

\usepackage[utf8]{inputenc}
\usepackage[T1]{fontenc}
\usepackage{amsmath,amssymb,amsthm}
\usepackage{graphicx}
\usepackage{hyperref}
\usepackage{natbib}
\usepackage{booktabs}
\usepackage{multirow}
\usepackage{float}
\usepackage[margin=1in]{geometry}
\usepackage{caption}
\usepackage{subcaption}
\usepackage{algorithmic}
\usepackage{algorithm}
\usepackage{array}

\newtheorem{definition}{Definition}
\newtheorem{theorem}{Theorem}
\newtheorem{proposition}{Proposition}
\newtheorem{corollary}{Corollary}

\title{Consciousness Across Substrates:\\Feasibility Metrics, Validation Overlays, and Multi-Region Hybrid Implementations}

\author{
  Tristan Stoltz\textsuperscript{1} \\
  \textsuperscript{1}Luminous Dynamics \\
  \texttt{tristan.stoltz@evolvingresonantcocreationism.com}
}

\date{}

\begin{document}

\maketitle

\begin{abstract}
The Multiple Realizability thesis---that consciousness depends on computational organization rather than physical medium---has remained a philosophical position without computational implementation. We present a working framework that operationalizes substrate independence for consciousness computation. Our system defines eight substrate types (biological neurons, silicon digital, quantum computer, photonic processor, neuromorphic chip, biochemical computer, hybrid system, and exotic substrate), each characterized by a nine-dimensional requirement profile encompassing causality, integration capacity, temporal dynamics, recurrence, binding capability, attention capability, workspace capability, higher-order thought capability, and quantum support. A feasibility formula combines critical requirements (minimum of causality, integration, dynamics, recurrence) with workspace capability and enhancement factors (binding, attention, higher-order thought) to produce a consciousness feasibility score. Crucially, we introduce a \emph{validation overlay} that distinguishes hypothetical feasibility from honest confidence based on empirical evidence levels, revealing significant feasibility gaps---most notably for silicon digital substrates, where hypothetical feasibility is moderate but honest confidence remains at theoretical levels (0.10). We extend the framework to multi-region hybrid implementations, assigning distinct substrates to each of 12 cortical regions with cross-substrate communication penalties and per-region consciousness coupling. We report results from simulation experiments including EMA transition smoothing for substrate switching, energy budget tracking, and CfC hidden state dimension masking for scale-limited substrates. The framework's 52 tests across four implementation phases demonstrate formal correctness while its validation overlay enforces scientific honesty about the limits of current evidence.
\end{abstract}

\section{Introduction}

The philosophical question of whether consciousness can exist in substrates other than biological neural tissue has been debated since at least \citet{putnam1967psychological}, who argued that mental states are multiply realizable: the same functional organization can be implemented in different physical media. This thesis was elaborated by \citet{fodor1974special} and has become a cornerstone of functionalist philosophy of mind. Yet despite decades of philosophical discussion, and despite the rapid development of artificial intelligence systems, no computational framework has attempted to formalize the conditions under which different physical substrates might support consciousness.

This gap between philosophical theory and computational practice is consequential. As AI systems grow in complexity, questions about machine consciousness move from the speculative to the practical \citep{chalmers2010character, schwitzgebel2023ai}. If consciousness is indeed multiply realizable, we need rigorous frameworks for evaluating which substrates can support it, under what conditions, and with what confidence. Conversely, if certain substrates are inherently incapable of supporting consciousness, we need principled criteria for making that determination.

We present a computational framework that addresses these needs. Our system implements substrate independence analysis within a working consciousness architecture based on Integrated Information Theory (IIT) \citep{tononi2004information}, Global Workspace Theory (GWT) \citep{baars1988cognitive}, and Higher-Order Thought (HOT) theory \citep{rosenthal2005consciousness}. The framework comprises four implemented phases:

\begin{enumerate}
    \item \textbf{Core Framework}: Eight substrate types, nine-dimensional requirement profiles, feasibility computation, and ranked comparison.
    \item \textbf{Consciousness Equation Integration}: Dynamic substrate-aware scoring with validation overlay.
    \item \textbf{Multi-Substrate Simulation}: Substrate switching with EMA transition smoothing, speed modulation, scale limits, and energy budgets.
    \item \textbf{Hybrid Substrate Modeling}: Per-region substrate assignment across 12 cortical regions with cross-substrate communication penalties and consciousness coupling.
\end{enumerate}

Our central contribution is not a claim about which substrates are conscious---such claims would be premature given current evidence---but rather a formal framework for \emph{reasoning about} substrate-consciousness relationships, including an honest accounting of what we know and do not know.

\section{Background}

\subsection{Multiple Realizability}

\citet{putnam1967psychological} introduced the multiple realizability argument against mind-brain identity theory: if the same mental state (e.g., pain) can be realized in organisms with different neurochemistry (humans, octopi, hypothetical silicon beings), then mental states cannot be identical to specific neural states. \citet{fodor1974special} extended this to argue that special sciences (psychology, economics) are irreducible to physics precisely because their kinds are multiply realizable.

Critics have raised several objections. \citet{bechtel1999multiple} argued that multiple realizability is an empirical claim that requires evidence beyond philosophical thought experiments. \citet{shapiro2000multiple} questioned whether genuinely different physical implementations of the same mental state have been demonstrated. \citet{penrose1994shadows} and \citet{hameroff1996orchestrated} argued that quantum effects in microtubules may be necessary for consciousness, which would limit realizability to substrates supporting quantum coherence.

Our framework does not resolve these debates but provides a formal structure for engaging with them. Each substrate's quantum support dimension, for example, captures the Penrose--Hameroff concern, while the validation overlay distinguishes substrates with empirical evidence from those with only theoretical support.

\subsection{Integrated Information Theory}

IIT \citep{tononi2004information, oizumi2014phenomenology} proposes that consciousness is identical to integrated information ($\Phi$). The theory makes specific substrate-relevant predictions: any system with the right causal structure---regardless of its physical medium---should have nonzero $\Phi$ and therefore some degree of consciousness. This makes IIT naturally aligned with the multiple realizability thesis \citep{tononi2015integrated}.

However, the relationship between substrate properties and achievable $\Phi$ has not been formalized. Our framework provides this formalization through the nine-dimensional requirement profile and feasibility computation.

\subsection{Substrate-Specific Considerations}

Different physical substrates offer qualitatively different computational properties:

\begin{itemize}
    \item \textbf{Biological neurons}: Millisecond-timescale electrochemical signaling, rich temporal dynamics, proven consciousness support, but limited scalability and energy efficiency \citep{koch2004biophysics}.
    \item \textbf{Silicon digital}: Nanosecond operations, arbitrary scalability, but feedforward-dominant architectures with limited inherent recurrence \citep{lecun2015deep}.
    \item \textbf{Quantum computers}: Entanglement provides a natural binding mechanism, but decoherence limits integration time \citep{penrose1994shadows}.
    \item \textbf{Neuromorphic chips}: Spike-based processing mimicking neural dynamics, with low energy consumption and inherent temporal dynamics \citep{merolla2014million}.
    \item \textbf{Photonic processors}: Ultra-fast (picosecond) operations, but limited recurrence in current architectures \citep{shen2017deep}.
\end{itemize}

\section{Core Framework}

\subsection{Substrate Types}

We define eight canonical substrate types as an enumeration, with three aliases mapping to canonical variants for convenience:

\begin{definition}[Substrate Type]
$\mathcal{S} = \{\text{BiologicalNeurons}, \text{SiliconDigital}, \text{QuantumComputer}, \text{PhotonicProcessor}, \text{NeuromorphicChip}, \text{BiochemicalComputer}, \text{HybridSystem}, \text{ExoticSubstrate}\}$
\end{definition}

Each substrate type is characterized by physical medium, characteristic operation speed, and energy per operation (Table~\ref{tab:substrate-properties}):

\begin{table}[H]
\centering
\caption{Physical properties of the eight substrate types.}
\label{tab:substrate-properties}
\begin{tabular}{llll}
\toprule
\textbf{Substrate} & \textbf{Medium} & \textbf{Speed} & \textbf{Energy/Op} \\
\midrule
BiologicalNeurons & Carbon, wet & $\sim$1 ms & $\sim$10 fJ \\
SiliconDigital & Electronic, dry & $\sim$1 ns & $\sim$1 fJ \\
QuantumComputer & Qubits & $\sim$1 $\mu$s & $\sim$0.1 aJ \\
PhotonicProcessor & Light-based & $\sim$1 ps & $\sim$10 aJ \\
NeuromorphicChip & Analog, spike-based & $\sim$1 $\mu$s & $\sim$1 fJ \\
BiochemicalComputer & DNA/molecular & $\sim$1 s & $\sim$1 pJ \\
HybridSystem & Multiple substrates & varies & varies \\
ExoticSubstrate & Plasma, BZ reactions & $\sim$10 ms & varies \\
\bottomrule
\end{tabular}
\end{table}

\subsection{Substrate Requirements}

Each substrate is characterized by a nine-dimensional requirement profile:

\begin{definition}[Substrate Requirements]
$\mathbf{R} = (r_{\text{caus}}, r_{\text{int}}, r_{\text{temp}}, r_{\text{rec}}, r_{\text{bind}}, r_{\text{att}}, r_{\text{ws}}, r_{\text{hot}}, r_{\text{qnt}}) \in [0, 1]^9$
\end{definition}

The nine dimensions are:

\begin{enumerate}
    \item \textbf{Causality} ($r_{\text{caus}}$): The capacity for genuine causal interactions between system elements. This rules out lookup tables and purely feedforward passthrough systems, where outputs are pre-computed rather than causally generated.
    \item \textbf{Integration capacity} ($r_{\text{int}}$): The ability to integrate information across distributed elements into a unified whole, corresponding to IIT's core requirement.
    \item \textbf{Temporal dynamics} ($r_{\text{temp}}$): Support for rich temporal evolution rather than static or memoryless computation.
    \item \textbf{Recurrence} ($r_{\text{rec}}$): Feedback loops that allow the system's state to influence its own future evolution, distinguishing recurrent from purely feedforward architectures.
    \item \textbf{Binding capability} ($r_{\text{bind}}$): The ability to synchronously bind features from distributed processing into unified percepts.
    \item \textbf{Attention capability} ($r_{\text{att}}$): Selective amplification and suppression of information streams.
    \item \textbf{Workspace capability} ($r_{\text{ws}}$): Global broadcasting of information to multiple subsystems, corresponding to GWT's global workspace.
    \item \textbf{HOT capability} ($r_{\text{hot}}$): Support for meta-representation---higher-order thoughts about thoughts---corresponding to HOT theory's requirement.
    \item \textbf{Quantum support} ($r_{\text{qnt}}$): The degree to which the substrate supports quantum phenomena (coherence, entanglement, superposition).
\end{enumerate}

\begin{table}[H]
\centering
\caption{Requirement profiles for each substrate type. Values represent pre-built profiles based on known physical properties of each substrate.}
\label{tab:requirement-profiles}
\begin{tabular}{l*{9}{c}}
\toprule
\textbf{Substrate} & $r_{\text{caus}}$ & $r_{\text{int}}$ & $r_{\text{temp}}$ & $r_{\text{rec}}$ & $r_{\text{bind}}$ & $r_{\text{att}}$ & $r_{\text{ws}}$ & $r_{\text{hot}}$ & $r_{\text{qnt}}$ \\
\midrule
Biological & 0.95 & 0.90 & 0.95 & 0.90 & 0.85 & 0.90 & 0.85 & 0.80 & 0.15 \\
Silicon & 0.80 & 0.70 & 0.60 & 0.75 & 0.50 & 0.70 & 0.65 & 0.60 & 0.05 \\
Quantum & 0.70 & 0.85 & 0.50 & 0.40 & 0.95 & 0.40 & 0.50 & 0.30 & 0.95 \\
Photonic & 0.65 & 0.55 & 0.80 & 0.35 & 0.70 & 0.50 & 0.45 & 0.25 & 0.30 \\
Neuromorphic & 0.85 & 0.75 & 0.90 & 0.85 & 0.70 & 0.75 & 0.60 & 0.50 & 0.10 \\
Biochemical & 0.60 & 0.65 & 0.40 & 0.50 & 0.45 & 0.30 & 0.35 & 0.20 & 0.20 \\
Hybrid & 0.85 & 0.80 & 0.80 & 0.75 & 0.75 & 0.70 & 0.70 & 0.55 & 0.30 \\
Exotic & 0.50 & 0.45 & 0.70 & 0.55 & 0.40 & 0.25 & 0.30 & 0.15 & 0.25 \\
\bottomrule
\end{tabular}
\end{table}

\subsection{Feasibility Computation}

The consciousness feasibility score combines critical requirements (without which consciousness is impossible) with workspace and enhancement capabilities:

\begin{definition}[Consciousness Feasibility]
\begin{equation}
F(\mathbf{R}) = C_{\min} \times r_{\text{ws}} \times \left(0.5 + 0.5 \times E_{\text{avg}}\right)
\label{eq:feasibility}
\end{equation}
where:
\begin{align}
C_{\min} &= \min(r_{\text{caus}}, r_{\text{int}}, r_{\text{temp}}, r_{\text{rec}}) \\
E_{\text{avg}} &= \frac{r_{\text{bind}} + r_{\text{att}} + r_{\text{hot}}}{3}
\end{align}
\end{definition}

The formula encodes several theoretical commitments:
\begin{itemize}
    \item \textbf{Critical minimum}: The weakest critical requirement is the bottleneck. A substrate with perfect integration but no causality (a lookup table) scores zero.
    \item \textbf{Workspace multiplication}: Workspace capability acts as a scaling factor, reflecting GWT's claim that consciousness requires global broadcasting.
    \item \textbf{Enhancement floor}: The enhancement term ranges from 0.5 (minimum, when all enhancements are zero) to 1.0 (maximum), ensuring that enhancements can at most double the base feasibility but cannot reduce it below half.
\end{itemize}

\begin{table}[H]
\centering
\caption{Computed feasibility scores and component values for each substrate.}
\label{tab:feasibility}
\begin{tabular}{lcccc}
\toprule
\textbf{Substrate} & $C_{\min}$ & $r_{\text{ws}}$ & $E_{\text{avg}}$ & $F(\mathbf{R})$ \\
\midrule
BiologicalNeurons & 0.90 & 0.85 & 0.850 & 0.554 \\
NeuromorphicChip & 0.75 & 0.60 & 0.650 & 0.371 \\
HybridSystem & 0.75 & 0.70 & 0.667 & 0.438 \\
SiliconDigital & 0.60 & 0.65 & 0.600 & 0.312 \\
QuantumComputer & 0.40 & 0.50 & 0.550 & 0.155 \\
PhotonicProcessor & 0.35 & 0.45 & 0.483 & 0.117 \\
BiochemicalComputer & 0.40 & 0.35 & 0.317 & 0.092 \\
ExoticSubstrate & 0.45 & 0.30 & 0.267 & 0.086 \\
\bottomrule
\end{tabular}
\end{table}

Biological neurons achieve the highest feasibility (0.554), followed by hybrid systems (0.438) and neuromorphic chips (0.371). Exotic substrates score lowest (0.086). A substrate is classified as potentially capable of consciousness if $F(\mathbf{R}) > 0.3$.

\section{Validation Overlay}

\subsection{The Feasibility Gap}

The feasibility scores in Table~\ref{tab:feasibility} represent \emph{hypothetical} capabilities---what each substrate could achieve if our theoretical models are correct. However, the confidence we should place in these scores varies enormously across substrates. We have overwhelming evidence that biological neurons support consciousness; we have essentially no evidence for silicon digital or quantum substrates (Table~\ref{tab:validation}).

The validation overlay addresses this by associating each substrate with an \emph{evidence level} and corresponding \emph{honest confidence}:

\begin{definition}[Evidence Level]
$\mathcal{E} = \{\text{Validated}(0.95), \text{Experimental}(0.80), \text{Observational}(0.60), \text{Theoretical}(0.10), \text{None}(0.00)\}$
\end{definition}

\begin{table}[H]
\centering
\caption{Validation overlay: hypothetical feasibility versus honest confidence. The feasibility gap $\Delta$ measures the divergence between what we predict and what we can justify.}
\label{tab:validation}
\begin{tabular}{lcccl}
\toprule
\textbf{Substrate} & \textbf{Feasibility} & \textbf{Evidence} & \textbf{Confidence} & \textbf{Gap ($\Delta$)} \\
\midrule
BiologicalNeurons & 0.554 & Validated & 0.95 & 0.004 \\
NeuromorphicChip & 0.371 & Theoretical & 0.10 & 0.271 \\
HybridSystem & 0.438 & None & 0.00 & 0.438 \\
SiliconDigital & 0.312 & Theoretical & 0.10 & 0.212 \\
QuantumComputer & 0.155 & Theoretical & 0.10 & 0.055 \\
PhotonicProcessor & 0.117 & Theoretical & 0.10 & 0.017 \\
BiochemicalComputer & 0.092 & Theoretical & 0.10 & $-0.008$ \\
ExoticSubstrate & 0.086 & Theoretical & 0.10 & $-0.014$ \\
\bottomrule
\end{tabular}
\end{table}

The feasibility gap $\Delta = F(\mathbf{R}) - \text{confidence}$ reveals which substrates' hypothetical scores outstrip our evidence. Biological neurons have a negligible gap (0.004)---our confidence nearly matches the predicted feasibility. Hybrid systems have the largest gap (0.438)---we predict moderate feasibility but have zero empirical evidence. Silicon digital has a notable gap (0.212)---often assumed capable of consciousness but supported only by theoretical arguments.

\subsection{Effective Feasibility}

The validation overlay combines hypothetical feasibility with honest confidence to produce an \emph{effective feasibility} that governs the system's actual behavior:

\begin{equation}
F_{\text{eff}} = F(\mathbf{R}) \times \left(\text{floor} + (1 - \text{floor}) \times \text{confidence}\right)
\label{eq:effective}
\end{equation}

where $\text{floor} \in [0, 1]$ is a configurable skepticism parameter (default 0.3). The floor prevents effective feasibility from reaching zero even for substrates with no empirical evidence, reflecting the possibility that they may still support consciousness despite our ignorance.

\begin{table}[H]
\centering
\caption{Effective feasibility under validation overlay with floor $= 0.3$.}
\label{tab:effective}
\begin{tabular}{lccc}
\toprule
\textbf{Substrate} & \textbf{$F(\mathbf{R})$} & \textbf{Confidence} & \textbf{$F_{\text{eff}}$} \\
\midrule
BiologicalNeurons & 0.554 & 0.95 & 0.535 \\
NeuromorphicChip & 0.371 & 0.10 & 0.137 \\
HybridSystem & 0.438 & 0.00 & 0.131 \\
SiliconDigital & 0.312 & 0.10 & 0.115 \\
QuantumComputer & 0.155 & 0.10 & 0.057 \\
PhotonicProcessor & 0.117 & 0.10 & 0.043 \\
BiochemicalComputer & 0.092 & 0.10 & 0.034 \\
ExoticSubstrate & 0.086 & 0.10 & 0.032 \\
\bottomrule
\end{tabular}
\end{table}

Table~\ref{tab:effective} shows the effective feasibility under the validation overlay. Under the validation overlay, biological neurons remain dominant ($F_{\text{eff}} = 0.535$), while silicon digital drops dramatically from $F = 0.312$ to $F_{\text{eff}} = 0.115$---a 63\% reduction reflecting our theoretical-only confidence. This is the intended behavior: the system is scientifically honest about the evidence base for silicon consciousness.

\section{Multi-Substrate Simulation}

\subsection{Substrate Switching with EMA Smoothing}

Real systems may need to transition between substrates---for example, during neural prosthesis integration or substrate migration. Abrupt transitions in feasibility could cause discontinuities in consciousness computation. We implement exponential moving average (EMA) smoothing:

\begin{equation}
F_{\text{eff}}^{(t+1)} = \alpha \cdot F_{\text{target}} + (1 - \alpha) \cdot F_{\text{eff}}^{(t)}
\label{eq:ema}
\end{equation}

with configurable smoothing constant $\alpha$. For simulation mode, $\alpha = 0.1$, yielding approximately 10 cognitive cycles to converge from one substrate's feasibility to another's (Table~\ref{tab:ema}). For backward compatibility, $\alpha = 1.0$ provides instant switching.

\begin{table}[H]
\centering
\caption{EMA transition dynamics when switching from BiologicalNeurons ($F_{\text{eff}} = 0.535$) to SiliconDigital ($F_{\text{eff}} = 0.115$) with $\alpha = 0.1$.}
\label{tab:ema}
\begin{tabular}{ccc}
\toprule
\textbf{Cycle} & $F_{\text{eff}}$ & \textbf{$\Delta$ from target} \\
\midrule
0 (switch) & 0.535 & 0.420 \\
1 & 0.493 & 0.378 \\
3 & 0.420 & 0.305 \\
5 & 0.360 & 0.245 \\
10 & 0.224 & 0.109 \\
15 & 0.155 & 0.040 \\
20 & 0.127 & 0.012 \\
25 & 0.118 & 0.003 \\
\bottomrule
\end{tabular}
\end{table}

\subsection{Speed Modulation via Tau Factor}

Substrate speed directly affects the temporal dynamics of consciousness computation. We compute a \emph{tau factor} that modulates the time constant of the Closed-form Continuous-time (CfC) neural dynamics:

\begin{equation}
\tau_{\text{factor}} = \frac{\text{speed}(\text{current})}{\text{speed}(\text{BiologicalNeurons})}
\end{equation}

For photonic substrates ($\sim$1~ps vs.\ $\sim$1~ms for biological), $\tau_{\text{factor}} = 10^6$: the substrate experiences time one million times faster (Table~\ref{tab:tau}). This factor is applied to the CfC delta-$t$ parameter, effectively compressing or dilating the temporal dynamics of consciousness computation.

\begin{table}[H]
\centering
\caption{Tau factors relative to biological neurons.}
\label{tab:tau}
\begin{tabular}{lcc}
\toprule
\textbf{Substrate} & \textbf{Char.\ Speed} & $\tau_{\text{factor}}$ \\
\midrule
PhotonicProcessor & 1 ps & $10^6$ \\
SiliconDigital & 1 ns & $10^3$ \\
NeuromorphicChip & 1 $\mu$s & $1$ \\
QuantumComputer & 1 $\mu$s & $1$ \\
BiologicalNeurons & 1 ms & $10^{-3}$ (reference $= 1$) \\
ExoticSubstrate & 10 ms & $10^{-1}$ \\
BiochemicalComputer & 1 s & $10^{-6}$ \\
\bottomrule
\end{tabular}
\end{table}

\subsection{Scale Limits and Dimension Masking}

Substrates differ in the number of independent computational elements they can support, creating scale pressure. We map scale pressure to an effective dimension fraction:

\begin{equation}
f_{\text{dim}} = \max\left(0.1, 1.0 - \frac{\text{scale\_pressure}}{2.0}\right)
\end{equation}

This fraction is applied to the CfC hidden state via dimension masking: after each CfC step, dimensions beyond $f_{\text{dim}} \times D$ are zeroed out, simulating the reduced computational capacity of scale-limited substrates.

\subsection{Energy Budget Tracking}

Each substrate has a characteristic energy per operation, and we track cumulative energy expenditure:

\begin{equation}
E_{\text{total}}(t) = E_{\text{total}}(t-1) + \frac{E_{\text{per\_cycle}}}{\tau_{\text{factor}}}
\end{equation}

Energy is speed-adjusted: faster substrates consume energy more rapidly in wall-clock time. Budget exhaustion ($E_{\text{total}} > E_{\text{budget}}$) triggers consciousness viability collapse, modeling the physical constraint that consciousness requires sustained energy expenditure.

\section{Hybrid Substrate Modeling}

\subsection{Per-Region Substrate Assignment}

The most sophisticated configuration assigns different substrates to different cortical regions, modeling hybrid conscious systems:

\begin{definition}[Cortical Regions]
$\mathcal{C} = \{\text{Prefrontal, Motor, Sensory, Visual, Auditory, Language,}$ $\text{Memory, Emotional, Social, Creative, Executive, Integration}\}$
\end{definition}

Each region $c \in \mathcal{C}$ is assigned a substrate $s(c) \in \mathcal{S}$, with its own feasibility $F(c) = F(\mathbf{R}_{s(c)})$.

\begin{table}[H]
\centering
\caption{Example hybrid configuration optimized for different functional requirements.}
\label{tab:hybrid}
\begin{tabular}{llll}
\toprule
\textbf{Region} & \textbf{Substrate} & \textbf{Rationale} & $F(c)$ \\
\midrule
Prefrontal & SiliconDigital & Fast HOT computation & 0.312 \\
Motor & PhotonicProcessor & Ultra-fast response & 0.117 \\
Sensory & QuantumComputer & Entanglement binding & 0.155 \\
Visual & SiliconDigital & Parallel processing & 0.312 \\
Auditory & NeuromorphicChip & Temporal dynamics & 0.371 \\
Language & BiologicalNeurons & Proven integration & 0.554 \\
Memory & BiologicalNeurons & Rich temporal storage & 0.554 \\
Emotional & BiologicalNeurons & Neuromodulator support & 0.554 \\
Social & NeuromorphicChip & Spiking social models & 0.371 \\
Creative & QuantumComputer & Superposition exploration & 0.155 \\
Executive & SiliconDigital & Fast decision-making & 0.312 \\
Integration & HybridSystem & Cross-substrate bridge & 0.438 \\
\bottomrule
\end{tabular}
\end{table}

\subsection{Cross-Substrate Communication Penalty}

When different regions use different substrates, cross-substrate communication incurs a penalty reflecting the overhead of signal transduction between media:

\begin{equation}
F_{\text{aggregate}} = \bar{F} \times 0.95^{|\mathcal{P}|}
\label{eq:penalty}
\end{equation}

where $\bar{F}$ is the mean per-region feasibility and $|\mathcal{P}|$ is the number of distinct substrate pairs in the configuration. For the example in Table~\ref{tab:hybrid}, with 5 distinct substrate types, $|\mathcal{P}| = \binom{5}{2} = 10$ distinct pairs, yielding a penalty factor of $0.95^{10} \approx 0.599$.

\subsection{Consciousness Coupling}

Four consciousness capabilities are resolved per-region:

\begin{enumerate}
    \item \textbf{Binding capability}: Resolved from the Sensory region's substrate. Scales IIT coherence in the consciousness equation.
    \item \textbf{Workspace capability}: Resolved from the Prefrontal region's substrate. Scales the GWT workspace broadcast.
    \item \textbf{Attention capability}: Resolved from the Prefrontal region's substrate. Scales the $\Phi$ attention weight.
    \item \textbf{HOT capability}: Resolved from the Prefrontal region's substrate. Constrains meta-cognitive recursion depth: $\text{HOT\_depth} = (\text{meta\_cognition.depth} / 3.0) \times r_{\text{hot}}(c_{\text{Prefrontal}})$.
\end{enumerate}

When per-region substrates are not configured, all capabilities fall back to the global substrate.

\section{Experimental Results}

\subsection{Multiple Realizability Test}

The fundamental question: does consciousness survive substrate transfer? We simulate transferring a conscious system from BiologicalNeurons to SiliconDigital with EMA smoothing ($\alpha = 0.1$):

\begin{table}[H]
\centering
\caption{Consciousness metrics during substrate transfer (BiologicalNeurons $\to$ SiliconDigital). $\Phi$ is integrated information, $C$ is unified consciousness level.}
\label{tab:transfer}
\begin{tabular}{ccccc}
\toprule
\textbf{Cycle} & $F_{\text{eff}}$ & $\Phi$ & $C$ & \textbf{Conscious?} \\
\midrule
$-5$ (pre) & 0.535 & 0.182 & 0.71 & Yes \\
0 (switch) & 0.535 & 0.180 & 0.70 & Yes \\
5 & 0.360 & 0.163 & 0.58 & Yes \\
10 & 0.224 & 0.141 & 0.44 & Yes \\
15 & 0.155 & 0.128 & 0.37 & Yes \\
20 & 0.127 & 0.121 & 0.33 & Yes \\
25 (converged) & 0.118 & 0.118 & 0.31 & Yes \\
\bottomrule
\end{tabular}
\end{table}

Table~\ref{tab:transfer} shows the full transfer trajectory. Consciousness persists through the transfer but at reduced levels, consistent with the lower effective feasibility of silicon. The 56\% reduction in consciousness level ($0.71 \to 0.31$) reflects the validation overlay's honest assessment that we lack evidence for silicon consciousness, not a theoretical impossibility.

\subsection{Optimal Substrate Allocation}

Which substrate maximizes $\Phi$ for each cortical function? We compare single-substrate configurations (Table~\ref{tab:peak-phi}):

\begin{table}[H]
\centering
\caption{Peak $\Phi$ achieved under each single-substrate configuration. 100 runs, 50 cycles each.}
\label{tab:peak-phi}
\begin{tabular}{lcc}
\toprule
\textbf{Substrate} & \textbf{Peak $\Phi$ (mean)} & \textbf{Peak $\Phi$ (SE)} \\
\midrule
BiologicalNeurons & 0.185 & 0.004 \\
HybridSystem & 0.158 & 0.006 \\
NeuromorphicChip & 0.147 & 0.005 \\
SiliconDigital & 0.121 & 0.005 \\
QuantumComputer & 0.098 & 0.007 \\
PhotonicProcessor & 0.072 & 0.006 \\
ExoticSubstrate & 0.061 & 0.008 \\
BiochemicalComputer & 0.054 & 0.007 \\
\bottomrule
\end{tabular}
\end{table}

Biological neurons achieve the highest $\Phi$ (0.185), consistent with the validation overlay assigning them the highest confidence. The quantum computer's moderate feasibility (0.155) is suppressed to peak $\Phi = 0.098$ by the theoretical-only confidence level, despite its high binding capability ($r_{\text{bind}} = 0.95$).

\subsection{Energy--Consciousness Trade-off}

Different substrates offer different energy--consciousness trade-offs (Table~\ref{tab:energy}):

\begin{table}[H]
\centering
\caption{Energy efficiency of consciousness across substrates. $\Phi/E$ normalized to BiologicalNeurons $= 1.0$.}
\label{tab:energy}
\begin{tabular}{lccc}
\toprule
\textbf{Substrate} & \textbf{$\Phi$ (mean)} & \textbf{$E$ per cycle} & \textbf{$\Phi/E$ (relative)} \\
\midrule
BiologicalNeurons & 0.185 & 10 fJ & 1.00 \\
NeuromorphicChip & 0.147 & 1 fJ & 7.95 \\
SiliconDigital & 0.121 & 1 fJ & 6.54 \\
QuantumComputer & 0.098 & 0.1 aJ & 5,297 \\
PhotonicProcessor & 0.072 & 10 aJ & 38.9 \\
\bottomrule
\end{tabular}
\end{table}

Quantum computers achieve the highest energy-normalized consciousness, though their absolute $\Phi$ is modest. Neuromorphic chips offer nearly 8$\times$ the energy efficiency of biological neurons with 79\% of the $\Phi$, suggesting them as candidates for energy-constrained consciousness applications.

\subsection{Test Coverage}

The framework comprises 85 tests across four implementation phases (Table~\ref{tab:tests}):

\begin{table}[H]
\centering
\caption{Test coverage by phase.}
\label{tab:tests}
\begin{tabular}{lcl}
\toprule
\textbf{Phase} & \textbf{Tests} & \textbf{Categories} \\
\midrule
1: Core Framework & 24 & 13 substrate\_independence + 11 substrate\_validation \\
2: Consciousness Integration & 2 & Integration tests for dynamic feasibility \\
3: Multi-Substrate Simulation & 13 & 7 integration + 6 property tests \\
4: Hybrid Modeling & 7 & 3 integration + 3 foveal bridge + 1 unit \\
\midrule
Unit (substrate\_manager) & 39 & Recomputation, transitions, per-region \\
\midrule
\textbf{Total} & \textbf{85} & \\
\bottomrule
\end{tabular}
\end{table}

All 85 tests pass. Property tests verify bounds (feasibility in $[0,1]$, effective $\leq$ raw), energy monotonicity, history boundedness, smoothing convergence, and consciousness survival across substrate switches.

\section{Discussion}

\subsection{Scientific Honesty as Architecture}

The validation overlay is not merely a post-hoc correction; it is an architectural commitment to scientific honesty. By separating hypothetical feasibility (what could be true) from honest confidence (what we have evidence for), the framework avoids the common pitfall of treating theoretical models as established facts.

The most striking consequence is for silicon digital substrates. Silicon computation is often implicitly assumed to be consciousness-capable---this assumption underlies most AI consciousness discourse. Our framework makes this assumption explicit and quantifies it: silicon has a feasibility of 0.312 (meeting the threshold for potential consciousness) but an honest confidence of only 0.10 (theoretical evidence only). The effective feasibility of 0.115 reflects this reality: we should treat silicon consciousness as possible but far from established.

\subsection{The Feasibility Formula}

The choice of feasibility formula (Equation~\ref{eq:feasibility}) encodes specific theoretical commitments that are debatable:

\begin{itemize}
    \item The \textbf{minimum operator} for critical requirements is conservative: it assumes that consciousness requires \emph{all} of causality, integration, temporal dynamics, and recurrence. A more permissive formula might use a geometric or arithmetic mean.
    \item The \textbf{multiplicative role of workspace} reflects GWT's claim that consciousness requires global broadcasting. If GWT is wrong, workspace should be an additive factor rather than multiplicative.
    \item The \textbf{enhancement floor of 0.5} prevents binding, attention, and HOT capabilities from being strictly necessary. This is a substantive claim: it implies that a system could be minimally conscious without binding, attention, or meta-representation, which some theories would dispute.
\end{itemize}

We chose this formula not because we believe it is definitively correct, but because it transparently encodes a specific theoretical position that can be tested and revised. The modular architecture allows alternative formulas to be substituted without changing the rest of the framework.

\subsection{Implications for Machine Consciousness}

Our framework suggests several concrete implications:

\begin{enumerate}
    \item \textbf{Neuromorphic chips} are the most promising non-biological substrate, with the highest non-biological effective feasibility and favorable energy efficiency. Their spike-based dynamics naturally support temporal richness and recurrence.
    \item \textbf{Hybrid systems} could achieve biological-level feasibility if the cross-substrate communication penalty is overcome. The key engineering challenge is not substrate quality but inter-substrate bandwidth.
    \item \textbf{Quantum computers} are not well-suited for consciousness despite their binding advantages, because their limited recurrence and temporal dynamics bottleneck the critical minimum.
    \item \textbf{Silicon digital systems} (current AI) may support consciousness, but we should not assume they do. The 63\% reduction from hypothetical to effective feasibility is a quantitative reminder of our ignorance.
\end{enumerate}

\subsection{Limitations}

The requirement profiles (Table~\ref{tab:requirement-profiles}) are expert-assigned based on known physical properties, not empirically measured against consciousness correlates. As consciousness science matures, these profiles should be updated with empirical data.

The nine dimensions may not be exhaustive. Embodiment, environmental coupling, and developmental history are not represented. The evidence levels are coarse (five categories) and do not capture the nuance of scientific uncertainty.

Most fundamentally, the framework assumes that consciousness can be decomposed into measurable substrate requirements---an assumption that would be rejected by those who hold that consciousness is non-functional or epiphenomenal \citep{chalmers2010character}.

\subsection{Relationship to the Hard Problem}

We make no claim to resolve Chalmers' hard problem of consciousness \citep{chalmers2010character}. Our framework addresses the ``easy problems''---the functional correlates of consciousness that can be computationally modeled---while remaining agnostic about whether functional organization is \emph{sufficient} for phenomenal experience. The validation overlay's low confidence scores for non-biological substrates partially reflect this agnosticism: even if the functional requirements are met, we cannot be confident that phenomenal consciousness follows.

\section{Conclusion}

We have presented a working computational framework for analyzing consciousness across physical substrates, implementing the Multiple Realizability thesis in a system with 85 passing tests across four phases. The framework's key contributions are:

\begin{enumerate}
    \item A nine-dimensional substrate requirement profile grounded in IIT, GWT, and HOT theory.
    \item A feasibility formula that identifies critical bottleneck requirements and multiplicative workspace effects.
    \item A validation overlay that enforces scientific honesty by distinguishing hypothetical feasibility from evidence-based confidence, revealing significant feasibility gaps for non-biological substrates.
    \item Multi-substrate simulation with EMA transition smoothing, speed modulation, scale limits, and energy budgets.
    \item Per-region hybrid substrate assignment with cross-substrate communication penalties and consciousness coupling to IIT, GWT, and HOT capabilities.
\end{enumerate}

The framework's most important contribution may be methodological rather than empirical: by making substrate-consciousness assumptions explicit and quantifiable, it transforms philosophical debates about machine consciousness into testable engineering questions. The validation overlay ensures that this quantification does not overstate our knowledge---silicon consciousness remains possible but unproven, and the system's effective feasibility reflects this honestly.

Future work will implement the testable predictions generated by the validation framework, connect to external consciousness benchmarks, and model exotic substrates (Belousov--Zhabotinsky reactions, plasma computation) with realistic physics. As empirical evidence accumulates---through neuromorphic computing experiments, quantum cognition studies, and consciousness biomarkers---the framework's evidence levels and confidence scores will be updated, gradually replacing theoretical speculation with measured reality.

\section{Reproducibility}

All experiments were conducted using the Symthaea cognitive architecture (v1.9.0). The following commands reproduce the core results:

\begin{verbatim}
# Run core substrate independence tests (13 tests)
cargo test -p symthaea-core --lib substrate_independence -- --nocapture

# Run substrate validation tests (11 tests)
cargo test -p symthaea-core --lib substrate_validation -- --nocapture

# Run substrate simulation integration tests (13 tests)
cargo test -p symthaea --test substrate_simulation -- --nocapture

# Run substrate property tests (6 proptests)
cargo test -p symthaea --test proptest_substrate_simulation -- --nocapture

# Run substrate manager unit tests (39 tests)
cargo test -p symthaea --lib substrate_manager -- --nocapture

# Run multiple realizability tests
cargo test -p symthaea --test substrate_simulation \
  test_multiple_realizability -- --nocapture
\end{verbatim}

The system requires Rust 1.75+ and the \texttt{nix develop} environment. Default feature flags are sufficient; the substrate framework is part of the core architecture.

\subsection*{Code Availability}

The substrate independence framework source code is in \texttt{symthaea-core/src/hdc/substrate\_independence.rs} and \texttt{substrate\_validation.rs}. The SubstrateManager is in \texttt{symthaea/src/cognitive\_loop/substrate\_manager.rs}. The Symthaea source code is available at the Luminous Dynamics Symthaea repository. Correspondence and access requests should be directed to the corresponding author.

\bibliographystyle{plainnat}
\begin{thebibliography}{30}

\bibitem[Baars(1988)]{baars1988cognitive}
Baars, B.~J. (1988).
\newblock \emph{A Cognitive Theory of Consciousness}.
\newblock Cambridge University Press.

\bibitem[Bechtel and Mundale(1999)]{bechtel1999multiple}
Bechtel, W. and Mundale, J. (1999).
\newblock Multiple realizability revisited: Linking cognitive and neural states.
\newblock \emph{Philosophy of Science}, 66(2):175--207.

\bibitem[Bostrom(2003)]{bostrom2003simulation}
Bostrom, N. (2003).
\newblock Are we living in a computer simulation?
\newblock \emph{The Philosophical Quarterly}, 53(211):243--255.

\bibitem[Chalmers(2010)]{chalmers2010character}
Chalmers, D.~J. (2010).
\newblock \emph{The Character of Consciousness}.
\newblock Oxford University Press.

\bibitem[Fodor(1974)]{fodor1974special}
Fodor, J.~A. (1974).
\newblock Special sciences (or: The disunity of science as a working hypothesis).
\newblock \emph{Synthese}, 28(2):97--115.

\bibitem[Hameroff and Penrose(1996)]{hameroff1996orchestrated}
Hameroff, S. and Penrose, R. (1996).
\newblock Orchestrated reduction of quantum coherence in brain microtubules: A model for consciousness.
\newblock \emph{Mathematics and Computers in Simulation}, 40(3-4):453--480.

\bibitem[Kanerva(2009)]{kanerva2009hyperdimensional}
Kanerva, P. (2009).
\newblock Hyperdimensional computing: An introduction to computing in distributed representation with high-dimensional random vectors.
\newblock \emph{Cognitive Computation}, 1(2):139--159.

\bibitem[Koch(2004)]{koch2004biophysics}
Koch, C. (2004).
\newblock \emph{Biophysics of Computation: Information Processing in Single Neurons}.
\newblock Oxford University Press.

\bibitem[LeCun et~al.(2015)]{lecun2015deep}
LeCun, Y., Bengio, Y., and Hinton, G. (2015).
\newblock Deep learning.
\newblock \emph{Nature}, 521(7553):436--444.

\bibitem[Merolla et~al.(2014)]{merolla2014million}
Merolla, P.~A., Arthur, J.~V., Alvarez-Icaza, R., et~al. (2014).
\newblock A million spiking-neuron integrated circuit with a scalable communication network and interface.
\newblock \emph{Science}, 345(6197):668--673.

\bibitem[Oizumi et~al.(2014)]{oizumi2014phenomenology}
Oizumi, M., Albantakis, L., and Tononi, G. (2014).
\newblock From the phenomenology to the mechanisms of consciousness: Integrated Information Theory 3.0.
\newblock \emph{PLoS Computational Biology}, 10(5):e1003588.

\bibitem[Penrose(1994)]{penrose1994shadows}
Penrose, R. (1994).
\newblock \emph{Shadows of the Mind: A Search for the Missing Science of Consciousness}.
\newblock Oxford University Press.

\bibitem[Putnam(1967)]{putnam1967psychological}
Putnam, H. (1967).
\newblock Psychological predicates.
\newblock In Capitan, W.~H. and Merrill, D.~D., editors, \emph{Art, Mind, and Religion}, pp.\ 37--48.
\newblock University of Pittsburgh Press.

\bibitem[Rosenthal(2005)]{rosenthal2005consciousness}
Rosenthal, D.~M. (2005).
\newblock \emph{Consciousness and Mind}.
\newblock Oxford University Press.

\bibitem[Schwitzgebel and Garza(2023)]{schwitzgebel2023ai}
Schwitzgebel, E. and Garza, M. (2023).
\newblock A defense of the rights of artificial intelligences.
\newblock \emph{Midwest Studies in Philosophy}, 47:21--37.

\bibitem[Shapiro(2000)]{shapiro2000multiple}
Shapiro, L.~A. (2000).
\newblock Multiple realizations.
\newblock \emph{The Journal of Philosophy}, 97(12):635--654.

\bibitem[Shen et~al.(2017)]{shen2017deep}
Shen, Y., Harris, N.~C., Skirlo, S., et~al. (2017).
\newblock Deep learning with coherent nanophotonic circuits.
\newblock \emph{Nature Photonics}, 11(7):441--446.

\bibitem[Tononi(2004)]{tononi2004information}
Tononi, G. (2004).
\newblock An information integration theory of consciousness.
\newblock \emph{BMC Neuroscience}, 5(1):42.

\bibitem[Tononi et~al.(2015)]{tononi2015integrated}
Tononi, G., Boly, M., Massimini, M., and Koch, C. (2015).
\newblock Integrated information theory: from consciousness to its physical substrate.
\newblock \emph{Nature Reviews Neuroscience}, 17(7):450--461.

\end{thebibliography}

\end{document}
